\documentclass[11pt]{article}
\usepackage[a4paper,margin=2.6cm]{geometry}
\usepackage{amsmath,amssymb,amsthm,booktabs}
\newtheorem{lemma}{Lemma}
\newtheorem{proposition}[lemma]{Proposition}
\newtheorem{corollary}[lemma]{Corollary}
\theoremstyle{remark}\newtheorem*{remark}{Remark}
\title{The discrete Landau count on a cosine window\\
\large kernel dimension of evaluation, and $\pi^2$ as the leading rung}
\author{D.\ Joubert}
\date{6 September 2026 --- companion to \texttt{the-well} v3, notebook \S133--135}
\begin{document}
\maketitle
\begin{abstract}
\noindent
On the $N{+}1$ cosine hats of $[0,L]$, the number of exponentially small
eigenvalues of the zero Gram is the maximal lead of the Nyquist count
over the node count,
$D_{\max}=\max_\gamma(\gamma L/2\pi-N_\Gamma(\gamma))$.
The lower bound is linear algebra, for an arbitrary node set: at every
spectral cut $\omega$ the evaluation map on hats of frequency ${<}\omega$
has kernel dimension at least $n(\omega)-N_\Gamma(\omega)$.
Those kernel vectors vanish on all nodes below $\omega$, so under RH
their Rayleigh quotients see only the leakage past $\omega$ and are
exponentially small.
The matching upper bound up to the $O(\log c)$ Slepian plunge is the
usual Landau--Widom width, not re-proved here; empirically
$\#\{\ell_k>2\}=\mathrm{round}(D_{\max})$.
The one-set Slepian inequality already in the repository, fed with a
set of measure $D_{\max}$ Nyquist cells, gives $\ell\le \pi^2 D_{\max}+\log(1/A)$.
The measured mean rung $11.0$ sits $10$--$15\%$ above $\pi^2\approx 9.87$;
that remainder is $\log(1/A)/D$ plus the hat-versus-PW discrepancy, not a
new constant. Nothing here bears on RH.
\end{abstract}

\section{The count}
Fix $L=\log\mu$ and the cosine window $V_N=\mathrm{span}\{\eta_n\}_{n=0}^{N}$
on $[0,L]$, frequencies $\omega_n=2\pi n/L$, band edge
$\omega_{\max}=2\pi N/L$. Let $\Gamma=\{\gamma_k\}_{k\ge 1}$ be any
increasing sequence of nodes (zeros, under RH; anything, for the Gram).
Write $N_\Gamma(\omega)=\#\{\gamma\in\Gamma:\gamma<\omega\}$ and
\[
n(\omega)=\#\{n=0,\ldots,N:\ \omega_n<\omega\}.
\]
The repository's
\[
D_{\max}
=\max\Bigl(\max_k\bigl(\gamma_k\tfrac{L}{2\pi}-k\bigr),\
N-N_\Gamma(\omega_{\max})\Bigr)
\]
is the continuous interpolation of $\max_\omega\bigl(n(\omega)-N_\Gamma(\omega)\bigr)$,
up to $O(1)$ (the constant hat $\eta_0$ is in every $V(\omega)$, so the
integer kernel lower bound is typically $D_{\max}+1$).

\begin{lemma}[evaluation kernel]
\label{lem:ker}
Let $\mathrm{Eval}_\omega:V(\omega)\to\mathbb C^{N_\Gamma(\omega)}$ send
$f\in V(\omega):=\mathrm{span}\{\eta_n:\omega_n<\omega\}$ to
$(\hat f(\gamma))_{\gamma<\omega}$. Then
\[
\dim\ker\mathrm{Eval}_\omega
\;\ge\;
n(\omega)-N_\Gamma(\omega).
\]
\end{lemma}
\begin{proof}
$V(\omega)$ has dimension $n(\omega)$. Point evaluations of entire
functions of exponential type are linear. The rank of $\mathrm{Eval}_\omega$
is at most the number of targets, $N_\Gamma(\omega)$. Rank-nullity gives
the claim. No property of $\Gamma$ is used.
\end{proof}

\begin{corollary}[arbitrary nodes]
\label{cor:D}
There is a subspace $K\subset V_N$ of dimension at least
$\max_\omega\bigl(n(\omega)-N_\Gamma(\omega)\bigr)$ consisting of
functions that vanish at every node in some $[0,\omega_\star]$.
\end{corollary}

That is (A) of \texttt{the-well}, as a lower bound, for any nodes on any
finite cosine window. The zeros of an $L$-function are a special case.

\section{Small eigenvalues}
Assume RH, so $Q_L(f)=\sum_\gamma\hat f(\gamma)^2$ on $W_L$. For
$f\in\ker\mathrm{Eval}_{\omega_\star}$,
\[
Q_L(f)
=\sum_{\gamma\ge\omega_\star}\hat f(\gamma)^2.
\]
A trigonometric polynomial of frequencies ${<}\omega_\star$ has Fourier
mass concentrated on $[0,\omega_\star]$ (Dirichlet kernel of the cosine
basis; the leakage past $\omega_\star$ is the tail of $\mathrm{sinc}$
sidelobes). Hence $Q_L(f)/\|f\|^2$ is exponentially small in the
time-bandwidth distance from $\omega_\star$ to the next nodes --- the
same desert/gap mechanism as \texttt{sampling-floor}, Proposition~2.
Those $\dim K$ directions are therefore exponentially small eigenvalues
of the Gram (and of $Q_N$ once the hat space resolves them).

The matching upper bound --- no more than $D_{\max}+O(\log c)$ small
eigenvalues --- is the Landau--Widom plunge: the concentration operator
on a set of measure $|E|$ has $O(\log(\tau|E|))$ eigenvalues in
$(e^{-C},1-e^{-C})$. Thresholding at $\ell>2$ nats (as in
\texttt{test\_depth\_law.py}) sits past that plunge on every window
measured. We do not redo Widom's analysis for the hat Gram; the
equality $\#\{\ell_k>2\}=\mathrm{round}(D_{\max})$ remains the
empirical closing of the $O(\log)$ gap.

\section{The leading rung is $\pi^2$}
The one-set Slepian bound already recorded in \texttt{FREEZE.md} and
\texttt{one-set-minorant.md} says that if $I^c$ samples $PW_\tau$,
\[
\ell
\;\le\;
\pi\,\tau\,|I|+\log(1/A)+o(\tau|I|),
\qquad \tau=L/2.
\]
Feed this with a set $I$ of $D_{\max}$ Nyquist cells,
$|I|=D_{\max}\cdot(2\pi/L)$. Then
\[
\pi\tau|I|
=\pi\cdot\frac L2\cdot D_{\max}\cdot\frac{2\pi}L
=\pi^2 D_{\max}.
\]
So
\begin{equation}
\label{eq:pi2}
\ell
\;\le\;
\pi^2 D_{\max}+\log(1/A)+o(D_{\max}),
\qquad
\pi^2\approx 9.87.
\end{equation}
The measured mean rung $\bar r=11.0$ is this leading term plus
$\log(1/A)/D_{\max}$. On the stored edge-value windows
(\texttt{report/discrete-landau.md}):

\begin{center}\small
\begin{tabular}{lrrrr}\toprule
window & $D_{\max}$ & $\ell/D$ & $\pi^2$ & $(\ell-\pi^2 D)/D$\\
\midrule
$\zeta$ 11 & 9.89 & 10.82 & 9.87 & $+0.95$\\
$\chi_3$ 16 / 38 & 4.96 / 12.32 & 11.21 / 11.32 & 9.87 & $+1.3$\\
$\chi_5$ 16 / 38 & 3.08 / 7.58 & 11.00 / 11.65 & 9.87 & $+1.1$ / $+1.8$\\
$\chi_8$ 16 & 1.77 & 11.22 & 9.87 & $+1.3$\\
$\chi_{29}$ 38 & 1.08 & 11.02 & 9.87 & $+1.1$\\
$\zeta$ 16$^\dagger$ & 15.00 & 9.31 & 9.87 & $-0.56$\\
$\chi_{13}$, $\chi_{31}$$^\ddagger$ & 1.19 / 0.90 & 9.13 / 9.47 & 9.87 & $-0.7$\\
\bottomrule
\end{tabular}\\[2pt]
{\footnotesize $\dagger$ sub-Nyquist at the edge: $D_{\max}$ still grows with $N$, $\ell$ saturates.
$\ddagger$ $D_{\max}\approx 1$: a single rung is not the mean rung.}
\end{center}

On the resolved degree-1 windows the remainder is $+1.0$ to $+1.8$ nats
per mode, i.e.\ $10$--$18\%$ above $\pi^2$. That is $\log(1/A)$ of size
$5$--$15$ nats spread over $D=3$--$12$, plus the cosine hats not being
exact prolates. It is not a second universal constant. Degree 2
($\bar r=8.7$--$9.4$ on 11a1) sits \emph{on} $\pi^2$, as the one-set
bound predicts once $A$ is not tiny.

The first rung $\sim 16$ and the last $\sim 5$ are the two ends of the
Slepian plunge of the longest piece (the desert): the most concentrated
mode pays $\sim 2c=L\gamma_1$ reduced by sampling, the last free mode
pays $O(1)$. Their average is the mean rung. The $N$-drift
$10.3\to 11.5$ from $N=41$ to $67$ is $\log(1/A)/D$ shrinking as $D$
grows, slowly, because $A$ itself depends on $N$ through $\omega_{\max}$.

\section{What is proved, what is not}
\begin{itemize}
\item \textbf{Proved.} Lemma~\ref{lem:ker} and Corollary~\ref{cor:D}:
the count from below, any nodes, no RH.
\item \textbf{Proved under RH, up to leakage estimates already in
\texttt{sampling-floor}.} Those kernel directions are exponentially
small eigenvalues of the Gram.
\item \textbf{Identified, not derived to the digit.} $\bar r=11$ is
$\pi^2$ plus a $1$~nat/mode remainder. Equation~\eqref{eq:pi2} is the
one-set bound of the repository applied to $D_{\max}$ cells; it is an
upper bound in the large-$c$ regime and sits $10$--$15\%$ low on the
windows we have.
\item \textbf{Not proved.} A threshold-free equality
$\#\{\lambda_k<e^{-C}\}=D_{\max}$ with an absolute $C$; a closed form
for $\log(1/A)$; transfer to unsaturated Galerkin windows
($\chi_3$ $\mu=80$: $D_{\max}$ counts dimensions $Q_N$ has not paid
for yet, \texttt{Dmax-vs-Q.md}).
\end{itemize}
Under RH, Lemma~\ref{lem:ker} plus the one-set bound \emph{is} the depth
law up to the remainder $\log(1/A)$. Without RH the same statements
hold for the Gram of whatever the nodes are. No prime enters.

\section*{Status}
Scripts: \texttt{code/dmax.py}, tests in \texttt{tests/test\_depth\_law.py}
(including a synthetic node set). Table: \texttt{report/discrete-landau.md}.
This note does not replace \texttt{the-well}; it answers (A) from below
and names the constant in (B).
\end{document}
